\SourceLeaf{304}
\Needspace{6\baselineskip}\begin{center}\textsc{Gauss} an \textsc{Schumacher}. Göttingen, 15. Mai 1843.\end{center}

Ich sende meinem letzten Briefe an Sie, mein theuerster Freund, noch ein Paar Worte
nach, um einem Buche eine Gerechtigkeit widerfahren zu lassen, welches ich bisher, zum
Theil in Folge eines vor zwanzig Jahren durch Sie (freilich sehr unschuldiger Weise)
veranlassten Umstandes, weniger beachtet habe, als es verdient.

Sie erinnern Sich vielleicht noch, dass Sie mir damals in Ihrem Hause eine an sich
ganz artige von \textsc{Möbius} aufgestellte geometrische Aufgabe vorlegten, wovon ich
die nichts zu wünschen lassende Auflösung fast augenblicklich niederschrieb. In Folge
dieses Vorfalls hatte ich das Buch von \textsc{Möbius}, barycentrischer Calcul, in
dessen Vorrede jene Aufgabe auch wieder erwähnt wurde, mit einer Art Vorurtheil in die
Hand genommen, nemlich mit dem Zweifel, ob es der Mühe werth sei, eine recht artig
ausgesonnene Rechnungsweise sich anzueignen, wenn man durch dieselbe nichts leisten
könne, was sich nicht eben so leicht ohne sie leisten lasse. Ich hatte deshalb das
Buch --- zumal da ich es in einer Zeit erhielt, wo andere Dinge mich ganz
beschäftigten ---, wie mir es freilich mit manchen Dingen gegangen ist, ohne viele
Erwartung davon zu haben, zunächst auf die Seite gelegt, und später völlig vergessen.

Eben in den allerletzten Tagen (richtiger gestern) fiel mir das Buch zufällig in die
Hände, und ich fand dann bald mit grossem Vergnügen, dass darin die Quintessenz der
Lehre von den Kegelschnitten in nucem gebracht ist, und dass gerade sein
barycentrischer Calcul auf dem leichtesten Wege zur Auflösung aller dahin gehörigen
Aufgaben führt. Interessant war mir darin auch manches andere beim Blättern sich
findende, z. B. der schöne von \textsc{Mac Laurin} und \textsc{Brackenridge}
gefundene, mir entweder unbekannt gebliebene, oder wieder vergessene Lehrsatz, mit
dessen Hülfe sich die Aufgabe, durch 5 gegebene Punkte einen Kegelschnitt zu legen, so
zierlich lösen lässt.
